\chapter{Introducción: por qué un FS basado en grafos}
\label{cap:intro}

Este libro es la presentación pedagógica del sistema sotFS: un
filesystem cuya estructura interna no es un árbol de inodes, como en
\texttt{ext4} o \texttt{XFS}, sino un \emph{grafo dirigido tipado} sobre
el que cada operación POSIX se ejecuta como una reescritura algebraica.
Antes de entrar en formalismos, conviene ver por qué la elección
estructural importa.

\section{Filesystems como árbol vs.\ como grafo}

\paragraph{El árbol POSIX clásico.} \texttt{ext4} y compañía modelan los
datos como un árbol jerárquico: directorios contienen inodes, inodes
apuntan a bloques. Es una decisión simple, eficiente y bien entendida.
Tres consecuencias se derivan de esa simplicidad:

\begin{itemize}
  \item Los \emph{hard links} no son una propiedad del modelo; son un
    bolt-on. Dos paths que apuntan al mismo inode existen, pero el modelo
    subyacente no los refleja como ``misma cosa con dos nombres''.
  \item La \emph{deduplicación} post-escritura es cara: hay que rescanear
    bloques iguales y reemplazar referencias.
  \item Los \emph{snapshots} se hacen vía copy-on-write o
    \emph{shadow paging}; son técnicas de implementación, no propiedades
    del modelo.
\end{itemize}

\paragraph{El grafo tipado.} sotFS usa una estructura distinta: un
\emph{grafo dirigido tipado} (TypeGraph) con seis tipos de nodo y siete
tipos de arista. Lo que en el árbol son convenciones aparece como
relaciones formales:

\begin{itemize}
  \item \textbf{Sharing entre dominios.} Un mismo bloque referenciado
    desde dos archivos sin bolt-on; emerge naturalmente del modelo.
  \item \textbf{Snapshots gratis.} Un snapshot es un puntero a un nodo
    raíz en cierto momento; nada se copia.
  \item \textbf{Capabilities first-class.} Una capability es un nodo del
    grafo, no una entrada en una tabla paralela.
  \item \textbf{Operaciones como reescritura.} Cada \opcreate, \opmkdir,
    \opunlink, etc.\ es una regla DPO sobre el grafo.
\end{itemize}

\paragraph{El precio.} Mayor complejidad teórica. Un grafo tipado tiene
invariantes que un árbol no necesita; las operaciones tienen
\emph{gluing conditions} formales; los chequeos son menos triviales. El
libro es básicamente la presentación ordenada de ese costo y de las
herramientas (TLA\textsuperscript{+}, Coq, fuzzers) para mantenerlo bajo
control.

\section{Lo que sotFS es y lo que no}

\paragraph{Lo que es.}
\begin{itemize}
  \item Un filesystem POSIX-compatible que monta en Linux vía \fuse{}.
  \item Un substrato algebraico (\TGraph) sobre el que se prueba la
    preservación de invariantes en TLA\textsuperscript{+} y Coq.
  \item Un workspace Rust de ocho crates (cap.~\ref{cap:implementacion}
    los recorre).
  \item Una herramienta de seguridad ofensiva-defensiva: provenance
    log con cap-mediated context, métricas estructurales (treewidth,
    curvatura), y functor de decepción $D_{\mathrm{FS}}$
    (caps.~\ref{cap:topologia} y \ref{cap:decepcion}).
\end{itemize}

\paragraph{Lo que no es (todavía).}
\begin{itemize}
  \item Un filesystem listo para producción a gran escala. La versión
    cubierta es \texttt{v0.2.5}; el roadmap planea Tier-2 (multi-host
    2PC), \texttt{sotfsctl repair}, y suite \texttt{xfstests/LTP}.
  \item Un sistema con prueba formal mecanizada \emph{end-to-end}. Hay
    cinco lemmas \texttt{Admitted} en Coq que el cap.~\ref{cap:coq}
    documenta uno por uno. La correspondencia DPO matemático $\to$ Rust
    es por inspección, no por extracción certificada.
  \item Un reemplazo del crate \texttt{sotfs-experimental}: ese crate
    es código aspiracional, no integrado al runtime estable, y queda
    \emph{fuera} de este libro.
\end{itemize}

\section{Mapa de capítulos}

La Tabla~\ref{tab:mapa} resume el contenido y las dependencias.

\begin{table}[h]
\centering
\renewcommand{\arraystretch}{1.15}
\begin{tabularx}{\linewidth}{c X l}
  \toprule
  \textbf{Cap.} & \textbf{Contenido} & \textbf{Depende de} \\
  \midrule
  \ref{cap:intro} & Introducción y motivación. & --- \\
  \ref{cap:prelim} & Preliminares: grafos tipados, pushouts, hashing. & 1 \\
  \ref{cap:typegraph} & El TypeGraph, los seis nodos y siete aristas, los siete invariantes. & 2 \\
  \ref{cap:hashing} & Hashing de contenido (\blake) y direccionamiento. & 3 \\
  \ref{cap:dpo} & Reescritura DPO de las operaciones POSIX. & 3, 4 \\
  \ref{cap:implementacion} & Implementación Rust: arenas, \rcu{}, motor de reglas. & 5 \\
  \ref{cap:persistencia} & Persistencia con \redb{} e índices secundarios. & 6 \\
  \ref{cap:capabilities} & Capabilities y delegación. & 3, 7 \\
  \ref{cap:transacciones} & GTXN y crash refinement. & 6, 7 \\
  \ref{cap:tla} & Verificación con TLA\textsuperscript{+}. & 5, 9 \\
  \ref{cap:coq} & Verificación con Coq, deuda Admitted explícita. & 5, 10 \\
  \ref{cap:topologia} & Treewidth y curvatura Ollivier-Ricci. & 3 \\
  \ref{cap:decepcion} & El functor de decepción $D_{\mathrm{FS}}$. & 3, 8 \\
  \ref{cap:fuse} & Adaptador \fuse{} y observabilidad. & 6, 8, 9 \\
  \bottomrule
\end{tabularx}
\caption{Mapa de capítulos y dependencias estrictas para una lectura
lineal. Los apéndices son lookup, no lectura lineal.}
\label{tab:mapa}
\end{table}

\section{Convenciones del libro}

\paragraph{Cajas tipadas.}
\begin{itemize}
  \item \emph{Definición} (violeta): introduce un objeto formal.
  \item \emph{Teorema}, \emph{Lema} (verde, teal): resultado central /
    auxiliar.
  \item \emph{Invariante} (azul): propiedad que el sistema mantiene
    siempre, a veces probada en TLA\textsuperscript{+} o Coq.
  \item \emph{Trampa} (rojo): error histórico del proyecto, con su
    \emph{lección} explícita; muchas enlazan a un commit (\citaCommit{<sha>}).
  \item \emph{Observación} (gris): nota lateral, no crítica.
  \item \emph{Ejercicio} (ámbar): propuesta al lector, graduada
    \dificultad{$\star$} / \dificultad{$\star\star$} / \dificultad{$\star\star\star$}.
\end{itemize}

\paragraph{Algoritmos.} En pseudocódigo estilo CLRS con
\texttt{algorithm2e}, numerados por capítulo. Cada algoritmo en el
libro tiene su contrapartida en código Rust real, citada por archivo y
línea: \citaLine{sotfs-graph/src/graph.rs}{571}.

\paragraph{Citas al repositorio.} Las macros \citaTLA{spec.tla} (specs
TLA\textsuperscript{+}), \citaCoq{archivo.v} (pruebas Coq),
\citaCrate{nombre} (crates Rust del workspace) y \citaCommit{sha}
(commits) puntúan referencias al código fuente. Las citas académicas
siguen el estilo numérico estándar y aparecen en la bibliografía al
final del volumen.

\section{Cómo leer este libro}

\paragraph{Lectura lineal.} Para alguien que quiere entender el sistema
de cero, leer en orden 1 a 14 es lo correcto. Cada capítulo descansa
sobre el anterior y los caps.~\ref{cap:typegraph} y \ref{cap:dpo} son
los más densos: si en algún momento te perdés en un capítulo posterior,
volvé a esos.

\paragraph{Lectura selectiva.}
\begin{description}
  \item[Implementadores en Rust:] caps.~\ref{cap:typegraph},
    \ref{cap:dpo}, \ref{cap:implementacion}, \ref{cap:persistencia} y
    los apéndices~\ref{ap:invariantes} (catálogo de invariantes) y
    \ref{ap:setup} (setup reproducible).
  \item[Verificacionistas:] caps.~\ref{cap:dpo}, \ref{cap:tla},
    \ref{cap:coq} y los apéndices~\ref{ap:dpo-tablas} (tablas DPO) y
    \ref{ap:tlc-results} (resultados TLC).
  \item[Practitioners de seguridad:] caps.~\ref{cap:capabilities},
    \ref{cap:topologia}, \ref{cap:decepcion}, \ref{cap:fuse}.
  \item[Educadores:] cualquier capítulo, con los ejercicios
    \dificultad{$\star$} y \dificultad{$\star\star$}; el
    apéndice~\ref{ap:soluciones} cubre las soluciones.
\end{description}

\section*{Resumen del capítulo}

\begin{itemize}
  \item sotFS abandona el árbol POSIX y modela su metadata como un
    grafo tipado con seis nodos y siete aristas.
  \item La elección habilita sharing, snapshots y capabilities como
    propiedades del modelo, no como bolt-ons.
  \item El precio se paga en complejidad teórica: invariantes,
    reescritura DPO, verificación.
  \item El libro recorre las catorce capas de ese costo, en orden
    pedagógico, con ejercicios graduados y apéndices de lookup.
\end{itemize}

\section*{Ejercicios}

\begin{ejercicio}[\dificultad{$\star$}]
\label{ej:intro-1}
Liste los seis tipos de nodo del TypeGraph mencionados en este capítulo.
(El cap.~\ref{cap:typegraph} los desarrolla.)
\end{ejercicio}

\begin{ejercicio}[\dificultad{$\star$}]
\label{ej:intro-2}
Tres consecuencias del modelo de árbol POSIX se mencionan como ``costos''.
Listelas en una línea cada una.
\end{ejercicio}

\begin{ejercicio}[\dificultad{$\star\star$}]
\label{ej:intro-3}
En el árbol POSIX, ¿cómo se implementan las hard links? ¿Cómo cambia el
mecanismo en sotFS según lo que adelanta este capítulo?
\end{ejercicio}

\begin{ejercicio}[\dificultad{$\star\star$}]
\label{ej:intro-4}
Dé un ejemplo concreto donde el sharing intrínseco de un Merkle DAG
ahorraría espacio respecto de \texttt{ext4}. ¿Cuántos bloques se
ahorrarían si dos archivos de 1\,MiB tienen 80\% de contenido en común?
\end{ejercicio}

\begin{ejercicio}[\dificultad{$\star\star\star$}]
\label{ej:intro-5}
Liste tres propiedades formales que sotFS apunta a probar y que
\texttt{ext4} \emph{no} prueba. Para cada una, anticipá qué herramienta
(TLA\textsuperscript{+}, Coq, o test/fuzz) sería la natural para
verificarla. (El cap.~\ref{cap:tla} y \ref{cap:coq} desarrollan esto.)
\end{ejercicio}
